\section{Problem Formulation}
\label{sec:problem-formulation}

\paragraph{Expert requests.} During generation, at routing step $t$ and
layer $\ell$, the model selects a set of experts
$R_t(\ell) \subseteq \{1,\dots,n_\text{experts}\}$ (the top-$k$ routed
experts for that layer and step). The sequence of per-layer routed sets
over all steps of a workload constitutes its \emph{expert-request
trace}. A step may batch multiple tokens (prompt processing), in which
case the experts required by every token in the step must be resident
\emph{simultaneously}; this report's replay model preserves step
boundaries rather than flattening them into independent per-token
accesses.

\paragraph{Byte-capacity retention.} Each expert at layer $\ell$ occupies
$\mathrm{bytes}(\ell)$ resident bytes, not assumed uniform across layers.
A \emph{retention policy} $\pi$ maintains, under a fixed byte budget
$\budget$, a resident set of (layer, expert) pairs, admitting and evicting
entries as routing proceeds, subject to
\begin{equation}
\sum_{(\ell,e)\,\in\,\text{resident}(\pi,t)} \mathrm{bytes}(\ell) \;\le\; \budget
\qquad \text{for every step } t. \label{eq:capacity}
\end{equation}
A \emph{miss} occurs when a requested expert is not resident and must be
fetched; a \emph{hit} occurs otherwise. We report both the miss
\emph{count} and the miss \emph{bytes} (the sum of $\mathrm{bytes}(\ell)$
over all missed requests); miss bytes are the primary quantity throughout
this report, since they are proportional to the volume of expert data
that would need to be serviced under the replay model, while miss count
alone does not account for per-layer size heterogeneity.

\paragraph{Causal vs.\ non-causal policies.} A policy is \emph{causal} if
its decision at step $t$ depends only on routing information available at
or before $t$ -- the only kind of policy that could ever be deployed in a
real runtime. A policy is \emph{non-causal} if it is permitted to use the
\emph{full} future routing sequence when making retention decisions at
step $t$. Every policy in this report is explicitly declared causal or
non-causal; interface-level assertions in the replay implementation
verify that no causal policy is accidentally given future information,
and that the sole non-causal policy is in fact built from full-trace
future information before replay begins.

\paragraph{The oracle.} The Belady-optimal offline policy
\cite{belady1966replacement} -- evict, at every eviction, whichever resident
item is needed again furthest in the future, or never again -- minimizes
the miss count of any policy operating under a fixed-capacity, single-cost
cache model, given full knowledge of the future reference sequence. We
use a global (cross-layer, cross-shard) variant of Belady,
\textsc{Oracle-Global-Belady}, as the non-causal upper-bound policy
throughout. It is unconditionally non-causal and is never treated as a
deployable baseline.

\begin{definition}[Information headroom]
\label{def:information-headroom}
For a causal baseline policy $\pi$ and the oracle $\pi^\star$, both
operating under the same byte budget $\budget$ and the same cache scope
(global, not shard-constrained), the \emph{relative information
headroom} at budget $\budget$ is
\begin{equation}
H_{\budget}(\pi) \;=\; \frac{\mathrm{missbytes}_\budget(\pi) -
\mathrm{missbytes}_\budget(\pi^\star)}{\mathrm{missbytes}_\budget(\pi)}.
\label{eq:information-headroom}
\end{equation}
$H_\budget(\pi)$ is an \emph{upper bound} on the fraction of $\pi$'s miss
bytes that any future causal policy -- no matter how well-informed by
observable routing structure -- could possibly avoid at that budget. It
is not itself a claim that any causal policy achieves this reduction.
\end{definition}

Section~\ref{sec:recoverable-envelope} formalizes the further
decomposition of this headroom away from the cost specifically
attributable to the real runtime's sharded retention architecture.
